\section{Procedura przekazania projektu}\label{sec:procedura}
\begin{enumerate}
\item Ustal plik główny, na przykład \texttt{main.tex}, i zapisz w
  \texttt{README.md} polecenie kompilacji.
\item Dodaj wszystkie pliki wczytywane przez \texttt{\textbackslash input},
  \texttt{\textbackslash includegraphics} i polecenia bibliografii.
\item Uruchom kompilację co najmniej dwa razy, a przy BibTeX użyj
  kolejności LaTeX -- BibTeX -- LaTeX -- LaTeX.
\item Otwórz PDF, sprawdź podpisy, numery odwołań, źródła i marginesy.
\item Pobierz lub skopiuj wyłącznie pliki źródłowe do nowego folderu,
  uruchom te same polecenia i porównaj wynik.
\end{enumerate}

Pliki pomocnicze takie jak \texttt{.aux}, \texttt{.log} czy
\texttt{.toc} nie muszą trafiać do repozytorium. Ich brak nie powinien
zmienić treści dokumentu, jeżeli komplet źródeł i polecenia są poprawne.
Gdy po skopiowaniu projekt przestaje się kompilować, sprawdź najpierw
ścieżki do grafik oraz pisownię nazw plików. Pomaga też przejrzenie
pierwszego istotnego błędu w pliku \texttt{.log}, zamiast poprawiania
wszystkich komunikatów po kolei.

Archiwizacja służy kontroli i dalszej edycji. Dla przykładu z
rozdziału~\ref{sec:cel} gotowy PDF ułatwia szybki podgląd, a pliki
źródłowe pozwalają zweryfikować, jak powstał.
